\documentclass[12pt]{article}
\usepackage[utf8]{inputenc}
\usepackage[T1]{fontenc}
\usepackage{hyperref}
\title{Tanzanian Sign Language Language Resources for K-12 Curriculum Development}
\author{Compiled by Claude Sonnet 5 (claude-sonnet-5) for Andr\'as Kornai}
\date{}
\begin{document}
\maketitle

\section{Introduction}
Tanzanian Sign Language (ISO 639-3 tza) is a deaf community sign language. The sources found describe several school-based varieties that developed from 1963. A standardized version was proposed in 1984 but has not been officially implemented.

\section{Primary Educational Resources}
\begin{itemize}
\item News report on the launch of the first digital sign language dictionary for students with hearing impairment, which also mentions a Curriculum Implementation Guide for secondary students with hearing disabilities: \url{https://www.thecitizen.co.tz/tanzania/oped/first-digital-dictionary-for-students-with-hearing-impairment-launched-2717340}
\item A second copy of the same report: \url{https://www.thecitizen.co.tz/oped/1840568-5629538-37vgju/index.html}
\end{itemize}

\section{Linguistic and Reference Materials}
\begin{itemize}
\item Wikipedia on Tanzanian sign languages: \url{https://en.wikipedia.org/wiki/Tanzanian_sign_languages}
\item University of Dodoma study on the challenges facing Tanzanian Sign Language learners in special schools for the deaf; it recommends teaching TSL as a subject and using it as a medium of instruction: \url{https://repository.udom.ac.tz/items/bf60274f-a301-4213-af02-eec11be38684}
\end{itemize}

\section{Religious and Community Resources}
None were found in the searches.

\section{Digital and Audio Resources}
The digital dictionary (Kiswahili, English and sign) is described in the news report above. I did not find or verify a direct link to the dictionary itself.
\begin{itemize}
\item Report on a sign language initiative launched by the Prime Minister: \url{https://inclusivenews.com.ng/tanzania-pm-launches-sign-language-initiative/}
\end{itemize}

\section{Curriculum Development Context}
Standardization is incomplete, and secondary and tertiary education are seldom offered in sign language. Any curriculum should work with the Tanzania Institute of Education and with deaf community groups to harmonize signs. The sources found are news reports and one academic study, so primary documentation is thin.
\end{document}